\section{Phase 7: Quantitative Convergence Theory for BLOM}
\label{sec:phase7_convergence}

The previous phases established BLOM as a well-posed local representation with exact locality, exact interpolation, and scheme-dependent global continuity.
What is still missing is a \emph{quantitative approximation theorem}: how close is BLOM-\(k\) to the exact global MINCO reference as the window size grows?

This section develops the first rigorous convergence theory for BLOM.
The main idea is to exploit the exponential off-diagonal decay of the inverse of the global banded optimality operator.
That decay, already identified in Phase~2 as the algebraic reason why exact localization is impossible, now becomes the source of quantitative approximation guarantees.

However, a distinction must be made between two objects:

\begin{enumerate}
	\item the \emph{idealized truncated map}, obtained by directly truncating the exact global sensitivity kernel;
	\item the \emph{actual BLOM-\(k\)} map, obtained from local window optimization and analytic reconstruction.
\end{enumerate}

The convergence theory completed in this chapter is fully rigorous for the first object, and becomes rigorous for the second object under an additional \emph{matching assumption}.
This limitation is made explicit below.

\subsection{Global exact coefficient map and centered sensitivity kernel}
\label{subsec:phase7_global_kernel}

We adopt the notation of Phases~1--2.
For fixed \(q,T\), the exact global MINCO coefficient vector is
\[
c^\star(q,T)\in\mathbb R^{2sM},
\]
and satisfies
\begin{equation}
	\label{eq:phase7_global_system}
	A(T)c^\star(q,T)=b(q,T),
\end{equation}
where \(A(T)\in\mathbb R^{2sM\times 2sM}\) is nonsingular and banded.

As in Phase~2, write the interior waypoint vector as
\[
\bar q:=(q_1,\dots,q_{M-1})^\top\in\mathbb R^{M-1},
\]
and decompose
\begin{equation}
	\label{eq:phase7_rhs_affine}
	b(q,T)=b_0(T)+S_q\,\bar q,
\end{equation}
where \(S_q\in\mathbb R^{2sM\times(M-1)}\) is the sparse selector matrix.
Define the centered exact coefficient map
\begin{equation}
	\label{eq:phase7_centered_c}
	\delta c^\star(\bar q,T)
	:=
	c^\star(\bar q,T)-c^\star(0,T).
\end{equation}
By Lemma~\ref{lem:phase2_exact_sensitivity} from Phase~2,
\begin{equation}
	\label{eq:phase7_centered_linear}
	\delta c^\star(\bar q,T)
	=
	A(T)^{-1}S_q\,\bar q.
\end{equation}

We now partition
\[
G(T):=A(T)^{-1}S_q \in \mathbb R^{2sM\times(M-1)}
\]
into block rows:
\begin{equation}
	\label{eq:phase7_G_blocks}
	G(T)=
	\begin{bmatrix}
		G_1(T)\\
		\vdots\\
		G_M(T)
	\end{bmatrix},
	\qquad
	G_i(T)\in\mathbb R^{2s\times(M-1)}.
\end{equation}
Let \(G_{ij}(T)\in\mathbb R^{2s\times 1}\) denote the \(j\)-th column block of \(G_i(T)\), so that
\begin{equation}
	\label{eq:phase7_ci_exact_kernel}
	\delta c_i^\star(\bar q,T)
	=
	\sum_{j=1}^{M-1} G_{ij}(T)\,q_j,
	\qquad i=1,\dots,M.
\end{equation}

Equation~\eqref{eq:phase7_ci_exact_kernel} is the exact global sensitivity expansion of the \(i\)-th coefficient block with respect to the sparse waypoint vector.

\subsection{Idealized truncated BLOM-\texorpdfstring{$k$}{k}}
\label{subsec:phase7_idealized_truncation}

The impossibility theorem of Phase~2 states that the exact map \eqref{eq:phase7_ci_exact_kernel} is not finitely local in general.
Nevertheless, if the kernel \(G_{ij}(T)\) decays exponentially in \(|i-j|\), then truncating the sum to a finite window yields an exponentially accurate local approximation.

\begin{definition}[Idealized truncated coefficient map]
	\label{def:phase7_idealized_truncation}
	For \(k\in\mathbb N\), define the index neighborhood
	\begin{equation}
		\label{eq:phase7_Nk}
		\mathcal N_k(i):=\{\,j\in\{1,\dots,M-1\}: |j-i|\le k\,\}.
	\end{equation}
	The idealized \(k\)-truncated coefficient of segment \(i\) is defined by
	\begin{equation}
		\label{eq:phase7_idealized_ck}
		\delta \tilde c_i^{(k)}(\bar q,T)
		:=
		\sum_{j\in\mathcal N_k(i)} G_{ij}(T)\,q_j.
	\end{equation}
	Equivalently,
	\begin{equation}
		\label{eq:phase7_idealized_ck_full}
		\tilde c_i^{(k)}(\bar q,T)
		=
		c_i^\star(0,T)+\delta \tilde c_i^{(k)}(\bar q,T).
	\end{equation}
\end{definition}

\begin{remark}
	\label{rem:phase7_idealized_truncation}
	The map \eqref{eq:phase7_idealized_ck} is the exact operator-theoretic BLOM-\(k\) surrogate suggested by Phase~2.
	It is not yet the same object as the actual BLOM-\(k\) produced by the local variational problems of Phases~3--4.
	Bridging the two requires an additional matching theorem; see Subsection~\ref{subsec:phase7_matching}.
\end{remark}

\subsection{Classical inverse-decay input}
\label{subsec:phase7_inverse_decay_input}

The main external tool of this section is a classical inverse-decay theorem for banded positive definite matrices.

\begin{theorem}[Classical inverse decay for banded positive definite matrices]
	\label{thm:phase7_demko}
	Let \(A\in\mathbb R^{n\times n}\) be symmetric positive definite and banded with scalar half-bandwidth \(\beta\).
	Assume that the spectrum of \(A\) is contained in \([\lambda_{\min},\lambda_{\max}]\) with
	\[
	0<\lambda_{\min}\le \lambda_{\max}<\infty.
	\]
	Then there exist constants
	\[
	C_0>0,
	\qquad
	\rho\in(0,1),
	\]
	depending only on \(\beta\) and the spectral condition number
	\[
	\kappa:=\frac{\lambda_{\max}}{\lambda_{\min}},
	\]
	such that
	\begin{equation}
		\label{eq:phase7_demko_scalar}
		|(A^{-1})_{ab}|
		\le
		C_0\,\rho^{|a-b|}
		\qquad
		\text{for all } a,b.
	\end{equation}
	By block embedding, the same conclusion holds at the block level: if \(A=[A_{ij}]_{i,j=1}^{M}\) is block-banded with fixed block size \(2s\), then
	\begin{equation}
		\label{eq:phase7_demko_block}
		\| (A^{-1})_{ij} \|_2
		\le
		C_1\,\rho^{|i-j|}
		\qquad
		\text{for all } i,j,
	\end{equation}
	for constants \(C_1>0\), \(\rho\in(0,1)\) depending only on the block bandwidth and spectral condition number.
\end{theorem}

\begin{remark}[Marked limitation]
	\label{rem:phase7_demko_limitation}
	Theorem~\ref{thm:phase7_demko} is used here as a \emph{classical external input}.
	Its proof is not reproduced in this chapter.
	What is proved below is how the inverse-decay estimate \eqref{eq:phase7_demko_block} implies quantitative BLOM error bounds.
\end{remark}

We next transfer the inverse-decay bound from \(A(T)^{-1}\) to the waypoint sensitivity kernel \(G(T)=A(T)^{-1}S_q\).

\begin{lemma}[Exponential decay of the waypoint sensitivity kernel]
	\label{lem:phase7_G_decay}
	Assume that the block inverse bound \eqref{eq:phase7_demko_block} holds for \(A(T)^{-1}\).
	Then there exist constants \(C_G>0\) and \(\rho\in(0,1)\) such that
	\begin{equation}
		\label{eq:phase7_G_decay}
		\|G_{ij}(T)\|_2
		\le
		C_G\,\rho^{|i-j|}
		\qquad
		\text{for all } i=1,\dots,M,\ j=1,\dots,M-1.
	\end{equation}
\end{lemma}

\begin{proof}
	Let \(s_j\in\mathbb R^{2sM}\) denote the \(j\)-th column of \(S_q\).
	By the Phase~1 construction, each interior waypoint \(q_j\) enters exactly one interpolation row block of the right-hand side.
	Therefore \(s_j\) is supported in exactly one block row, say \(r(j)\), and
	\[
	\|s_j\|_2 = 1.
	\]
	Moreover, the block-row index \(r(j)\) differs from the knot index \(j\) by at most a fixed offset \(\delta_q\) independent of \(M\) (indeed, in the canonical ordering one may take \(\delta_q=0\) or \(1\), depending on the precise block convention).
	Now
	\[
	G_{ij}(T)
	=
	\bigl(A(T)^{-1}S_q\bigr)_{ij}
	=
	(A(T)^{-1})_{i,r(j)}\,s_j.
	\]
	Hence
	\[
	\|G_{ij}(T)\|_2
	\le
	\|(A(T)^{-1})_{i,r(j)}\|_2\,\|s_j\|_2
	\le
	C_1\,\rho^{|i-r(j)|}.
	\]
	Using \(|i-r(j)|\ge |i-j|-\delta_q\), we obtain
	\[
	\|G_{ij}(T)\|_2
	\le
	C_1\,\rho^{-\delta_q}\rho^{|i-j|}
	=
	C_G\,\rho^{|i-j|},
	\]
	where \(C_G:=C_1\rho^{-\delta_q}\).
\end{proof}

\subsection{Level 1: uniform time allocation}
\label{subsec:phase7_uniform_time}

We first consider the easiest regime:
\[
T_i\equiv h>0.
\]
In this case, the exact global matrix becomes a banded matrix family depending only on the common step length \(h\) and the problem size \(M\).
To obtain an error bound uniform in \(M\), one needs uniform spectral control of that family.

\begin{assumption}[Uniform-time spectral stability]
	\label{ass:phase7_uniform_spectral}
	Fix \(s\) and \(h>0\).
	For every \(M\), let \(A_h^{(M)}\) denote the exact global coefficient matrix of the Phase~1 mother problem under
	\[
	T_1=\cdots=T_M=h.
	\]
	Assume that there exist constants
	\[
	0<\underline\lambda_h\le \overline\lambda_h<\infty
	\]
	independent of \(M\) such that
	\begin{equation}
		\label{eq:phase7_uniform_spectral}
		\underline\lambda_h I \preceq A_h^{(M)} \preceq \overline\lambda_h I
		\qquad
		\text{for all } M.
	\end{equation}
\end{assumption}

\begin{remark}[Marked limitation]
	\label{rem:phase7_uniform_spectral_limitation}
	A complete proof of Assumption~\ref{ass:phase7_uniform_spectral} from first principles --- for general \(s\) and with an explicit closed-form decay factor \(\rho_h\) --- is \textbf{not completed in this chapter}.
	What is completed is the implication
	\[
	\text{uniform spectral stability}
	\Longrightarrow
	\text{exponential BLOM-\(k\) approximation}.
	\]
\end{remark}

\begin{theorem}[Uniform-time coefficient error bound for idealized BLOM-\(k\)]
	\label{thm:phase7_uniform_error}
	Assume Assumption~\ref{ass:phase7_uniform_spectral}.
	Then there exist constants
	\[
	C_h>0,\qquad \rho_h\in(0,1),
	\]
	depending only on \(s\) and \(h\), such that for every \(M\), every interior segment \(i\), every \(k\in\mathbb N\), and every interior waypoint vector \(\bar q\),
	\begin{equation}
		\label{eq:phase7_uniform_error}
		\bigl\|
		\delta \tilde c_i^{(k)}(\bar q,h)-\delta c_i^\star(\bar q,h)
		\bigr\|_2
		\le
		C_h\,\rho_h^{\,k}\,\|\bar q\|_\infty .
	\end{equation}
\end{theorem}

\begin{proof}
	By Assumption~\ref{ass:phase7_uniform_spectral}, the condition numbers of the family \(A_h^{(M)}\) are uniformly bounded:
	\[
	\kappa(A_h^{(M)})\le \frac{\overline\lambda_h}{\underline\lambda_h}=: \kappa_h.
	\]
	Since the block bandwidth is fixed once \(s\) is fixed, Theorem~\ref{thm:phase7_demko} yields constants \(C_G>0\) and \(\rho_h\in(0,1)\), independent of \(M\), such that
	\[
	\|G_{ij}(h)\|_2\le C_G \rho_h^{|i-j|}.
	\]
	Using \eqref{eq:phase7_ci_exact_kernel} and \eqref{eq:phase7_idealized_ck},
	\begin{align}
		\delta c_i^\star-\delta \tilde c_i^{(k)}
		&=
		\sum_{j\notin\mathcal N_k(i)} G_{ij}(h)\,q_j.
		\label{eq:phase7_tail_start}
	\end{align}
	Taking norms and using \(\|q_j\|=|q_j|\le \|\bar q\|_\infty\),
	\begin{align}
		\bigl\|
		\delta c_i^\star-\delta \tilde c_i^{(k)}
		\bigr\|_2
		&\le
		\sum_{j\notin\mathcal N_k(i)}
		\|G_{ij}(h)\|_2\,|q_j|
		\notag\\
		&\le
		C_G\|\bar q\|_\infty
		\sum_{j\notin\mathcal N_k(i)}
		\rho_h^{|i-j|}.
		\label{eq:phase7_tail_bound_1}
	\end{align}
	For each distance \(r\ge k+1\), there are at most two indices \(j\) such that \(|i-j|=r\).
	Hence
	\begin{align}
		\sum_{j\notin\mathcal N_k(i)} \rho_h^{|i-j|}
		&\le
		2\sum_{r=k+1}^{\infty}\rho_h^{r}
		=
		2\frac{\rho_h^{k+1}}{1-\rho_h}.
		\label{eq:phase7_geometric_tail}
	\end{align}
	Substituting \eqref{eq:phase7_geometric_tail} into \eqref{eq:phase7_tail_bound_1} gives
	\[
	\bigl\|
	\delta c_i^\star-\delta \tilde c_i^{(k)}
	\bigr\|_2
	\le
	\frac{2C_G\rho_h}{1-\rho_h}\,\rho_h^k\,\|\bar q\|_\infty.
	\]
	Setting
	\[
	C_h:=\frac{2C_G\rho_h}{1-\rho_h}
	\]
	proves \eqref{eq:phase7_uniform_error}.
\end{proof}

\subsection{Level 2: bounded nonuniform times}
\label{subsec:phase7_bounded_nonuniform}

We next generalize the previous result to bounded nonuniform time allocations:
\[
0<T_{\min}\le T_i\le T_{\max}.
\]
As before, the key issue is whether the exact global matrix family remains uniformly well-conditioned as \(M\) varies.

\begin{assumption}[Uniform spectral stability on a bounded time box]
	\label{ass:phase7_box_spectral}
	Fix \(s\), \(T_{\min}>0\), and \(T_{\max}>T_{\min}\).
	Assume that there exist constants
	\[
	0<\underline\lambda_\ast\le \overline\lambda_\ast<\infty
	\]
	such that for every \(M\) and every
	\[
	T\in[T_{\min},T_{\max}]^M,
	\]
	the corresponding Phase~1 global coefficient matrix satisfies
	\begin{equation}
		\label{eq:phase7_box_spectral}
		\underline\lambda_\ast I \preceq A(T) \preceq \overline\lambda_\ast I.
	\end{equation}
\end{assumption}

\begin{remark}[Marked limitation]
	\label{rem:phase7_box_spectral_limitation}
	A fully general proof of Assumption~\ref{ass:phase7_box_spectral}, uniform in \(M\), is \textbf{not completed in this chapter}.
	For each fixed \(M\), continuity of eigenvalues on the compact box \([T_{\min},T_{\max}]^M\) gives existence of positive lower and upper spectral bounds.
	What remains open here is the proof that such bounds can be chosen uniformly over the entire family in \(M\).
\end{remark}

\begin{theorem}[Bounded-nonuniform coefficient error bound for idealized BLOM-\(k\)]
	\label{thm:phase7_box_error}
	Assume Assumption~\ref{ass:phase7_box_spectral}.
	Then there exist constants
	\[
	C_\ast>0,\qquad \rho_\ast\in(0,1),
	\]
	depending only on \(s\), \(T_{\min}\), and \(T_{\max}\), such that for every \(M\), every admissible \(T\in[T_{\min},T_{\max}]^M\), every interior segment \(i\), every \(k\in\mathbb N\), and every \(\bar q\),
	\begin{equation}
		\label{eq:phase7_box_error}
		\bigl\|
		\delta \tilde c_i^{(k)}(\bar q,T)-\delta c_i^\star(\bar q,T)
		\bigr\|_2
		\le
		C_\ast\,\rho_\ast^{\,k}\,\|\bar q\|_\infty.
	\end{equation}
\end{theorem}

\begin{proof}
	The proof is identical to that of Theorem~\ref{thm:phase7_uniform_error}.
	By Assumption~\ref{ass:phase7_box_spectral}, all admissible matrices \(A(T)\) have condition numbers uniformly bounded by
	\[
	\kappa_\ast:=\frac{\overline\lambda_\ast}{\underline\lambda_\ast}.
	\]
	Hence Theorem~\ref{thm:phase7_demko} and Lemma~\ref{lem:phase7_G_decay} provide constants \(C_G^\ast>0\) and \(\rho_\ast\in(0,1)\) such that
	\[
	\|G_{ij}(T)\|_2\le C_G^\ast \rho_\ast^{|i-j|}
	\]
	uniformly for all admissible \(T\).
	The same geometric-tail argument as in \eqref{eq:phase7_tail_start}--\eqref{eq:phase7_geometric_tail} then yields \eqref{eq:phase7_box_error}.
\end{proof}

\subsection{From idealized truncation to actual BLOM-\texorpdfstring{$k$}{k}}
\label{subsec:phase7_matching}

Theorems~\ref{thm:phase7_uniform_error} and \ref{thm:phase7_box_error} are rigorous statements for the idealized truncated map \(\tilde c_i^{(k)}\).
To transfer them to the actual BLOM-\(k\) obtained from local window optimization, one needs a comparison theorem between the two constructions.

\begin{assumption}[Matching of actual BLOM-\(k\) to the idealized truncation]
	\label{ass:phase7_matching}
	Assume that for the actual BLOM-\(k\) coefficient map \(c_i^{\mathrm{BLOM},k}\), there exist constants
	\[
	C_{\mathrm{match}}>0,\qquad \rho_{\mathrm{match}}\in(0,1),
	\]
	such that
	\begin{equation}
		\label{eq:phase7_matching}
		\bigl\|
		c_i^{\mathrm{BLOM},k}(\bar q,T)-\tilde c_i^{(k)}(\bar q,T)
		\bigr\|_2
		\le
		C_{\mathrm{match}}\,\rho_{\mathrm{match}}^{\,k}\,\|\bar q\|_\infty
	\end{equation}
	for all admissible data.
\end{assumption}

\begin{remark}[Marked limitation]
	\label{rem:phase7_matching_limitation}
	A proof of Assumption~\ref{ass:phase7_matching} for the fully general actual BLOM-\(k\) constructed from the windowwise local variational problems of Phases~3--4 is \textbf{not completed in this chapter}.
	This is the principal missing step between the idealized truncation theory and a full convergence theorem for the implemented BLOM-\(k\).
\end{remark}

\begin{corollary}[Conditional convergence of actual BLOM-\(k\)]
	\label{cor:phase7_actual_convergence}
	Assume either the hypotheses of Theorem~\ref{thm:phase7_uniform_error} or those of Theorem~\ref{thm:phase7_box_error}, together with Assumption~\ref{ass:phase7_matching}.
	Then there exist constants \(C_{\mathrm{tot}}>0\) and \(\rho_{\mathrm{tot}}\in(0,1)\) such that
	\begin{equation}
		\label{eq:phase7_actual_error}
		\bigl\|
		c_i^{\mathrm{BLOM},k}(\bar q,T)-c_i^\star(\bar q,T)
		\bigr\|_2
		\le
		C_{\mathrm{tot}}\,\rho_{\mathrm{tot}}^{\,k}\,\|\bar q\|_\infty.
	\end{equation}
\end{corollary}

\begin{proof}
	By the triangle inequality,
	\begin{align*}
		\bigl\|
		c_i^{\mathrm{BLOM},k}-c_i^\star
		\bigr\|_2
		&\le
		\bigl\|
		c_i^{\mathrm{BLOM},k}-\tilde c_i^{(k)}
		\bigr\|_2
		+
		\bigl\|
		\tilde c_i^{(k)}-c_i^\star
		\bigr\|_2.
	\end{align*}
	The first term is bounded by \eqref{eq:phase7_matching}; the second is bounded by either \eqref{eq:phase7_uniform_error} or \eqref{eq:phase7_box_error}.
	Thus
	\[
	\bigl\|
	c_i^{\mathrm{BLOM},k}-c_i^\star
	\bigr\|_2
	\le
	C_1\rho_1^k\|\bar q\|_\infty + C_2\rho_2^k\|\bar q\|_\infty.
	\]
	Let
	\[
	\rho_{\mathrm{tot}}:=\max(\rho_1,\rho_2)\in(0,1),
	\qquad
	C_{\mathrm{tot}}:=C_1+C_2.
	\]
	Then
	\[
	\bigl\|
	c_i^{\mathrm{BLOM},k}-c_i^\star
	\bigr\|_2
	\le
	C_{\mathrm{tot}}\,\rho_{\mathrm{tot}}^{\,k}\,\|\bar q\|_\infty.
	\]
\end{proof}

\subsection{Level 3: suboptimality of the control cost}
\label{subsec:phase7_cost_bound}

We now turn coefficient error bounds into cost bounds.
This requires additional care: the exact MINCO coefficient vector \(c^\star\) is the minimizer of the Phase~1 mother problem over its global affine feasible set.
Hence the quadratic cost orthogonality formula is valid only if the BLOM-\(k\) candidate is itself \emph{globally feasible for that same affine set}.

For fixed \((q,T)\), let
\[
\mathcal F(q,T)\subset \mathbb R^{2sM}
\]
denote the affine feasible coefficient set of the Phase~1 global mother problem.
Choose an orthonormal basis matrix \(N(T)\) of the corresponding homogeneous tangent subspace, so every feasible coefficient vector can be written uniquely as
\begin{equation}
	\label{eq:phase7_feasible_param}
	c=c_p(q,T)+N(T)z,
	\qquad
	z\in\mathbb R^{d_T},
\end{equation}
for some feasible anchor \(c_p(q,T)\).
Let the global minimum-control cost be
\begin{equation}
	\label{eq:phase7_global_cost}
	J(c,T)=\frac12\,c^\top H(T)c,
\end{equation}
where \(H(T)\succeq 0\) is the segmentwise snap Gram matrix from Phase~1.
Then on the feasible affine set,
\begin{equation}
	\label{eq:phase7_reduced_cost}
	\widehat J(z;q,T)
	=
	J(c_p+Nz,T)
	=
	\frac12 z^\top R(T) z + r(q,T)^\top z + \gamma(q,T),
\end{equation}
where
\begin{equation}
	\label{eq:phase7_reduced_H}
	R(T):=N(T)^\top H(T)N(T)\succeq 0.
\end{equation}
Because the global minimum is unique, \(R(T)\) is in fact positive definite on the reduced feasible coordinates.

Let \(z^\star\) be the exact MINCO reduced coordinate and \(z^{(k)}\) any other feasible reduced coordinate.
Set
\[
e_z:=z^{(k)}-z^\star.
\]

\begin{lemma}[Exact quadratic cost gap on the feasible affine set]
	\label{lem:phase7_exact_cost_gap}
	If both \(c^\star\) and \(c^{(k)}\) belong to the same feasible affine set \(\mathcal F(q,T)\), then
	\begin{equation}
		\label{eq:phase7_exact_cost_gap}
		J(c^{(k)},T)-J(c^\star,T)
		=
		\frac12\, e_z^\top R(T)e_z.
	\end{equation}
\end{lemma}

\begin{proof}
	Since \(z^\star\) minimizes \(\widehat J(z;q,T)\), it satisfies the first-order optimality condition
	\[
	R(T)z^\star + r(q,T)=0.
	\]
	Now expand \(\widehat J(z^\star+e_z;q,T)\):
	\begin{align*}
		\widehat J(z^\star+e_z;q,T)
		&=
		\frac12(z^\star+e_z)^\top R(z^\star+e_z)
		+
		r^\top(z^\star+e_z)
		+
		\gamma \\
		&=
		\left(\frac12 z^{\star\top}Rz^\star+r^\top z^\star+\gamma\right)
		+
		e_z^\top(Rz^\star+r)
		+
		\frac12 e_z^\top R e_z.
	\end{align*}
	The middle term vanishes by optimality:
	\[
	Rz^\star+r=0.
	\]
	Hence
	\[
	\widehat J(z^\star+e_z;q,T)-\widehat J(z^\star;q,T)
	=
	\frac12 e_z^\top R e_z.
	\]
	Since \(\widehat J\) is exactly the restriction of \(J\) to the feasible affine set, this proves \eqref{eq:phase7_exact_cost_gap}.
\end{proof}

\begin{theorem}[Absolute cost gap bound]
	\label{thm:phase7_absolute_cost_gap}
	Assume that a BLOM-\(k\) coefficient vector \(c^{\mathrm{BLOM},k}(q,T)\) is globally feasible for the same affine set \(\mathcal F(q,T)\) as the Phase~1 mother problem.
	Assume further that its coefficient error satisfies
	\begin{equation}
		\label{eq:phase7_error_for_cost}
		\bigl\|
		c^{\mathrm{BLOM},k}(q,T)-c^\star(q,T)
		\bigr\|_2
		\le
		C_e\,\rho^k\,\|\bar q\|_\infty.
	\end{equation}
	Then
	\begin{equation}
		\label{eq:phase7_absolute_cost_bound}
		0
		\le
		J^{\mathrm{BLOM},k}(q,T)-J^{\mathrm{MINCO}}(q,T)
		\le
		\frac12\,\lambda_{\max}(R(T))\,C_e^2\,\rho^{2k}\,\|\bar q\|_\infty^2.
	\end{equation}
\end{theorem}

\begin{proof}
	Because \(N(T)\) is chosen with orthonormal columns, the coefficient error inside the feasible affine set satisfies
	\[
	c^{\mathrm{BLOM},k}-c^\star = N(T)e_z,
	\]
	and therefore
	\[
	\|e_z\|_2 = \|c^{\mathrm{BLOM},k}-c^\star\|_2.
	\]
	By Lemma~\ref{lem:phase7_exact_cost_gap},
	\[
	J^{\mathrm{BLOM},k}-J^{\mathrm{MINCO}}
	=
	\frac12 e_z^\top R(T)e_z
	\le
	\frac12 \lambda_{\max}(R(T))\|e_z\|_2^2.
	\]
	Substituting \eqref{eq:phase7_error_for_cost} gives
	\[
	J^{\mathrm{BLOM},k}-J^{\mathrm{MINCO}}
	\le
	\frac12 \lambda_{\max}(R(T))\,C_e^2\,\rho^{2k}\,\|\bar q\|_\infty^2.
	\]
	Nonnegativity follows because \(c^\star\) is the global minimizer over \(\mathcal F(q,T)\).
\end{proof}

\begin{corollary}[Relative cost bound]
	\label{cor:phase7_relative_cost}
	Under the assumptions of Theorem~\ref{thm:phase7_absolute_cost_gap}, assume in addition that on the admissible set of interest there exists a constant \(\underline\alpha>0\) such that
	\begin{equation}
		\label{eq:phase7_minco_lower_bound}
		J^{\mathrm{MINCO}}(q,T)\ge \underline\alpha\,\|\bar q\|_\infty^2.
	\end{equation}
	Then
	\begin{equation}
		\label{eq:phase7_relative_cost_bound}
		\frac{J^{\mathrm{BLOM},k}(q,T)}{J^{\mathrm{MINCO}}(q,T)}
		\le
		1+
		\frac{\lambda_{\max}(R(T))\,C_e^2}{2\underline\alpha}\,\rho^{2k}.
	\end{equation}
	In particular,
	\[
	\frac{J^{\mathrm{BLOM},k}}{J^{\mathrm{MINCO}}}
	=
	1+O(\rho^{2k}).
	\]
\end{corollary}

\begin{proof}
	Divide \eqref{eq:phase7_absolute_cost_bound} by \(J^{\mathrm{MINCO}}(q,T)\), then use the lower bound \eqref{eq:phase7_minco_lower_bound}.
\end{proof}

\begin{remark}[Marked limitation for the cost theorem]
	\label{rem:phase7_cost_limitation}
	Theorem~\ref{thm:phase7_absolute_cost_gap} and Corollary~\ref{cor:phase7_relative_cost} are \textbf{conditional} on global feasibility with respect to the Phase~1 mother problem.
	At the current stage, the Phase~5 assembly schemes A/B/C do not yet come with a theorem guaranteeing that their assembled trajectories belong to the exact global feasible set of Phase~1 (which involves stronger global compatibility than the currently proved \(C^{s-1}\) continuity).
	Therefore, the relative cost bound is rigorously proved here only as a \emph{conditional theorem}.
	A full application of it to the present assembled BLOM trajectories remains open.
\end{remark}

\subsection{Interpretation of the three levels}
\label{subsec:phase7_interpretation}

The three levels of convergence theory established in this chapter should be interpreted as follows.

\paragraph{Level 1 (uniform times).}
If the global coefficient matrix family is spectrally stable under
\[
T_i\equiv h,
\]
then the idealized BLOM truncation converges exponentially with rate \(\rho_h^k\).
The proof is rigorous modulo the classical inverse-decay theorem and the spectral stability assumption.
An explicit closed-form formula for \(\rho_h\) is \textbf{not completed here}.

\paragraph{Level 2 (bounded nonuniform times).}
If the global coefficient matrix family remains uniformly spectrally stable on the time box
\[
[T_{\min},T_{\max}]^M,
\]
then the same exponential convergence holds uniformly:
\[
\bigl\|
\delta \tilde c_i^{(k)}-\delta c_i^\star
\bigr\|
\le
C_\ast \rho_\ast^k \|\bar q\|_\infty.
\]
Again, this theorem is rigorous under the explicit assumption of family-wise spectral stability.

\paragraph{Level 3 (cost suboptimality).}
If, in addition, actual BLOM-\(k\) is both:
\begin{enumerate}
	\item exponentially close to the idealized truncation, and
	\item globally feasible for the Phase~1 mother problem,
\end{enumerate}
then its control effort is relatively optimal up to an \(O(\rho^{2k})\) multiplicative factor.

\subsection{What has and has not been proved in Phase~7}
\label{subsec:phase7_status}

We summarize the exact status of the present convergence theory.

\begin{enumerate}
	\item \textbf{Fully proved in this chapter:}
	\begin{itemize}
		\item the idealized truncated BLOM-\(k\) error bound under exponential kernel decay;
		\item the transfer from inverse decay to coefficient truncation error;
		\item the conditional cost gap and relative cost bound on the common feasible affine set.
	\end{itemize}
	
	\item \textbf{Proved only under explicit assumptions:}
	\begin{itemize}
		\item uniform-time convergence, under family-wise spectral stability;
		\item bounded-nonuniform convergence, under uniform spectral stability on the admissible time box;
		\item convergence of the actual BLOM-\(k\), under the matching assumption to the idealized truncation.
	\end{itemize}
	
	\item \textbf{Not completed in this chapter:}
	\begin{itemize}
		\item a general proof that the family \(A(T)\) is uniformly spectrally stable over all \(M\);
		\item an explicit closed-form formula for the decay factor \(\rho\) in the canonical minimum-snap case;
		\item a general proof that the actual BLOM-\(k\) constructed in Phases~3--4 satisfies the matching estimate \eqref{eq:phase7_matching};
		\item a proof that the currently assembled BLOM trajectories of Phase~5 are globally feasible for the exact Phase~1 mother problem.
	\end{itemize}
\end{enumerate}

Accordingly, the mathematically correct summary of Phase~7 is:
\[
\boxed{
	\text{BLOM admits a rigorous exponential approximation theory at the idealized truncated-kernel level,}
}
\]
and
\[
\boxed{
	\text{the same rate transfers to actual BLOM-\(k\) and to cost suboptimality once the matching and feasibility steps are proved.}
}
\]

\section{Unfinished Proofs: Their Impact, Priority, and Technical Bottlenecks}
\label{sec:unfinished_proofs_discussion}

This section summarizes the current status of the main proofs that are not yet fully completed in the BLOM theory.
The purpose is not to conceal the gaps, but to state clearly:
\begin{enumerate}
	\item what each missing proof is supposed to establish,
	\item how its absence affects the overall theory,
	\item whether it is worth pursuing immediately,
	\item and where the present mathematical bottleneck lies.
\end{enumerate}

The discussion below is intentionally written in plain language, but kept in an academically precise form.

\subsection{Overview: not all missing proofs have the same importance}
\label{subsec:unfinished_overview}

At the current stage, the unfinished proofs do \emph{not} all affect the theory at the same level.
Their importance is stratified as follows:

\begin{itemize}
	\item The most important missing proof is the \emph{matching theorem} between actual BLOM-\(k\) and the idealized truncated BLOM-\(k\).
	Without it, the convergence theory proved in Phase~7 applies rigorously only to an idealized object, not yet to the exact algorithmic object implemented in BLOM.
	
	\item The second most important missing proof is the \emph{global feasibility theorem} showing that the current Phase~5 assembled BLOM trajectories automatically belong to the Phase~1 feasible set.
	Without it, the coefficient approximation theory can still stand, but the control-cost suboptimality theorem remains conditional.
	
	\item The proof of \emph{family-wise spectral stability for arbitrary \(M\)} is important for turning the convergence theory into a uniform theorem with constants independent of trajectory length.
	However, its absence does not invalidate the local-approximation idea itself.
	
	\item The derivation of an \emph{explicit closed-form decay factor \(\rho\)} in the canonical minimum-snap case is the least urgent.
	Its absence affects the sharpness and interpretability of the theory, but not the main qualitative convergence conclusion.
\end{itemize}

Thus, the correct strategic reading is:

\begin{quote}
	The unfinished proofs do not all threaten the same part of the theory.
	Some block the transfer from ideal theory to the actual BLOM construction; others only prevent the theory from being fully general or aesthetically complete.
\end{quote}

\subsection{Family-wise spectral stability for general \texorpdfstring{$M$}{M}}
\label{subsec:unfinished_spectral_stability}

\paragraph{What this proof is supposed to show.}
The goal is to prove that the global MINCO coefficient matrix family
\[
A(T)\in\mathbb R^{2sM\times 2sM}
\]
remains uniformly well conditioned as the number of segments \(M\) varies, provided the duration vector \(T\) stays in a bounded admissible set.
More concretely, one wants uniform constants
\[
0<\underline\lambda\le \lambda_{\min}(A(T))\le \lambda_{\max}(A(T))\le \overline\lambda<\infty
\]
that do not deteriorate as \(M\to\infty\).

\paragraph{Why this matters.}
If such a theorem were available, then the inverse-decay constants in Phase~7 could be chosen uniformly over all problem sizes.
This would upgrade the convergence theorem from a ``fixed-size or assumption-based result'' to a true family-wise theorem.

\paragraph{What happens if this proof is missing.}
Without this theorem, one can still prove convergence:
\begin{itemize}
	\item for fixed \(M\),
	\item or under an explicit spectral-stability assumption,
	\item or with strong numerical evidence over a large but finite range of \(M\).
\end{itemize}
Therefore, the theory does \emph{not} collapse.
What is lost is the clean statement that a single pair of constants \((C,\rho)\) works uniformly for all sufficiently large problem sizes.

\paragraph{Should this be pursued immediately?}
Not necessarily.
At the present stage, this proof is not the one that most directly decides whether BLOM is mathematically meaningful.
It is better viewed as a \emph{generalization and strengthening theorem}, rather than the first theorem that must be closed.

\paragraph{Where the bottleneck is.}
For each fixed \(M\), compactness and continuity arguments are enough to control the spectrum.
The real difficulty is the passage to ``uniformly in \(M\)''.
That step requires additional structure, such as:
\begin{itemize}
	\item block Toeplitz or asymptotically Toeplitz analysis,
	\item discrete coercivity estimates,
	\item uniform Poincar\'e-type inequalities,
	\item or symbol-based spectral theory.
\end{itemize}
In other words, the obstruction is not algebraic bookkeeping, but the lack of a genuinely \emph{family-level} spectral argument.

\subsection{Matching theorem: actual BLOM-\texorpdfstring{$k$}{k} versus idealized truncated BLOM-\texorpdfstring{$k$}{k}}
\label{subsec:unfinished_matching}

\paragraph{What this proof is supposed to show.}
Phase~7 proves exponential convergence for the idealized truncated object
\[
\tilde c_i^{(k)},
\]
which is obtained by truncating the exact global sensitivity kernel
\[
A(T)^{-1}S_q.
\]
The missing matching theorem would show that the actual BLOM-\(k\) coefficient
\[
c_i^{\mathrm{BLOM},k}
\]
satisfies
\[
\|c_i^{\mathrm{BLOM},k}-\tilde c_i^{(k)}\|\le C\rho^k\|\bar q\|,
\]
or at least an estimate of the same order.

\paragraph{Why this matters.}
This is the single most important unfinished proof.
Without it, the rigorous Phase~7 convergence theory applies to the \emph{idealized truncated kernel model}, but not yet fully to the actual BLOM map generated from:
\begin{itemize}
	\item local variational windows,
	\item natural boundary conditions,
	\item analytic local solves,
	\item and segment extraction / assembly.
\end{itemize}

\paragraph{What happens if this proof is missing.}
Then the theory is forced into the following weaker but honest form:
\begin{quote}
	We have a rigorous convergence theorem for an idealized BLOM-like truncation, and strong numerical evidence that actual BLOM-\(k\) appears to behave the same way.
\end{quote}
This is still meaningful, but it leaves a gap between the \emph{theoretical object} and the \emph{algorithmic object}.

\paragraph{Should this be pursued immediately?}
Yes.
Among all unfinished proofs, this is the one most worth attacking first.
It directly determines whether the theory established in Phase~7 is genuinely a theorem \emph{about BLOM}, rather than about a nearby idealization.

\paragraph{Where the bottleneck is.}
The bottleneck is conceptual:
the actual BLOM-\(k\) construction is not defined by truncating the global inverse kernel.
It is defined by:
\begin{itemize}
	\item solving a local variational problem,
	\item imposing natural boundary conditions,
	\item reconstructing a local polynomial,
	\item and extracting the central segment.
\end{itemize}
The idealized truncated map, by contrast, is defined directly by cutting off the exact global linear sensitivity operator.
The two constructions are intuitively close, but they are not literally the same formula.

Therefore, proving the matching theorem likely requires a new bridge, for example:
\begin{itemize}
	\item a Schur-complement interpretation of the local window solve,
	\item a Green's-function localization identity,
	\item or a consistency theorem relating the local Euler--Lagrange solution to the truncated global influence kernel.
\end{itemize}
At present, that bridge has not yet been built.

\subsection{Automatic feasibility of the Phase~5 assembled BLOM trajectories in the Phase~1 feasible set}
\label{subsec:unfinished_feasibility}

\paragraph{What this proof is supposed to show.}
The cost comparison theorem of Phase~7 is sharp only if both the exact MINCO solution and the BLOM-\(k\) candidate lie in the same affine feasible set of the Phase~1 mother problem.
Thus, the missing theorem would show that the assembled BLOM trajectory produced in Phase~5 automatically belongs to the exact global feasible set used in Phase~1.

\paragraph{Why this matters.}
This proof is essential if one wants to claim not only coefficient closeness, but also rigorous control-cost near-optimality:
\[
\frac{J^{\mathrm{BLOM}-k}}{J^{\mathrm{MINCO}}}\le 1+O(\rho^{2k}).
\]

\paragraph{What happens if this proof is missing.}
Without it, the coefficient approximation theory can still stand.
One may still prove statements of the form
\[
\|c^{(k)}-c^\star\|\le C\rho^k\|\bar q\|.
\]
However, one can no longer directly invoke the exact quadratic orthogonality identity behind the global cost-gap theorem.
Thus:
\begin{itemize}
	\item the coefficient theory survives,
	\item but the strongest cost-theoretic theorem remains conditional.
\end{itemize}

\paragraph{Should this be pursued immediately?}
It depends on the intended claim of the paper.
If the paper is positioned primarily as:
\begin{quote}
	a local, differentiable, sparse, MINCO-compatible trajectory representation,
\end{quote}
then this theorem can be postponed.
If the paper is positioned as:
\begin{quote}
	a quantitatively near-optimal replacement for global MINCO in control effort,
\end{quote}
then this theorem becomes high priority.

\paragraph{Where the bottleneck is.}
The bottleneck is structural, not technical.
The current Phase~5 assembly schemes prove:
\begin{itemize}
	\item Scheme A/B: global \(C^{s-1}\),
	\item Scheme C: exact \(C^0\), and conditional higher continuity.
\end{itemize}
But the exact Phase~1 feasible set is tied to the global MINCO mother problem, which has stronger compatibility structure than the presently proved assembly properties.
So the issue is not simply that the proof is long; it is that the currently assembled BLOM object and the exact Phase~1 feasible object are not yet manifestly the same mathematical object.
At this point, one must either:
\begin{enumerate}
	\item strengthen the assembly mechanism, or
	\item weaken the comparison theorem to a relaxed feasible set.
\end{enumerate}

\subsection{Explicit formula for the decay factor \texorpdfstring{$\rho$}{rho} in the canonical minimum-snap case}
\label{subsec:unfinished_explicit_rho}

\paragraph{What this proof is supposed to show.}
The current convergence theorems assert the existence of a decay factor
\[
\rho\in(0,1)
\]
such that the error behaves like \(O(\rho^k)\).
A stronger theorem would derive an explicit or semi-explicit formula for \(\rho\) in the canonical minimum-snap case.

\paragraph{Why this matters.}
An explicit \(\rho\) would improve:
\begin{itemize}
	\item interpretability,
	\item predictive power,
	\item and the ability to estimate how much window growth improves accuracy.
\end{itemize}
It would make the convergence theory much more informative.

\paragraph{What happens if this proof is missing.}
The qualitative theory still stands.
One can still state:
\[
\|c_i^{(k)}-c_i^\star\|\le C\rho^k\|\bar q\|
\]
for some \(\rho\in(0,1)\).
Thus, the missing explicit formula affects the \emph{sharpness} and \emph{readability} of the theorem, not its essential validity.

\paragraph{Should this be pursued immediately?}
No.
Among the unfinished proofs, this is the least urgent.
It is a strengthening result, not a foundational one.

\paragraph{Where the bottleneck is.}
The matrix family in the canonical minimum-snap problem is not a simple scalar tridiagonal Toeplitz matrix.
It is:
\begin{itemize}
	\item block-banded,
	\item induced by Hermite and Gram structures,
	\item and perturbed by boundary effects.
\end{itemize}
The inverse-decay theorem guarantees existence of some \(\rho\), but extracting a closed-form \(\rho\) requires substantially finer spectral analysis:
\begin{itemize}
	\item block symbol analysis,
	\item infinite periodic extension,
	\item or explicit spectral radius bounds.
\end{itemize}
This is mathematically deep, but not the first obstacle blocking the main BLOM theory.

\subsection{Priority ranking}
\label{subsec:unfinished_priority}

The unfinished proofs should be prioritized as follows.

\begin{enumerate}
	\item \textbf{Highest priority: matching theorem.}
	This determines whether the convergence theorem is truly about actual BLOM-\(k\).
	
	\item \textbf{Second priority: feasibility theorem.}
	This determines whether the strongest control-cost near-optimality theorem can be stated unconditionally.
	
	\item \textbf{Third priority: family-wise spectral stability.}
	This upgrades the theory from a conditional or fixed-size result to a uniform family theorem.
	
	\item \textbf{Lowest priority: explicit \(\rho\) formula.}
	This improves sharpness and interpretability, but is not needed to make the main theory stand.
\end{enumerate}

\subsection{Practical interpretation}
\label{subsec:unfinished_practical_interpretation}

The most practical way to interpret the current status is the following:

\begin{quote}
	The BLOM theory is already strong enough to justify continued development, numerical validation, and implementation.
	What remains incomplete are primarily the bridge theorems that connect the ideal operator-theoretic picture to the exact algorithmic BLOM construction, and the strongest cost-optimality claims.
\end{quote}

Therefore, the unfinished proofs do \emph{not} imply that the project should stop.
They imply only that the strongest final theorems should not yet be declared as closed.
A good immediate strategy is:
\begin{enumerate}
	\item continue numerical validation of Phase~7,
	\item use the results to decide whether actual BLOM behaves like the ideal truncation,
	\item and then decide whether to invest further effort into the matching theorem and the feasibility theorem.
\end{enumerate}

\subsection{Summary}
\label{subsec:unfinished_summary}

In one sentence:

\begin{quote}
	The missing matching theorem is the main unresolved bridge between the proven idealized convergence theory and the actual BLOM algorithm; the missing feasibility theorem is the main unresolved bridge between coefficient approximation and rigorous control-cost near-optimality; the remaining two missing proofs mainly affect generality and sharpness rather than the basic validity of the BLOM program.
\end{quote}